% Pursuing Stacks (cotes 134-2 to 134-8): archive cote 134-5, site page 123 (PDF page 124), typescript page 310.
% typescript. Draft compared: ps5.tex:899-964.
% AI-transcribed by Claude Opus 5.5 (claude-opus-5-5, adaptive thinking, effort high), Message Batches API, batch ps3 (msgbatch_01XdFmuVVBAcgc7EFxdtCqsj), 2026-11-12; not yet checked by eye.
\scanpage{134-5/123}
\mspage{310}
\[i_n = i_{\hand{!}}\mathrm{Sd}^n\hand{$\alpha$} : \hand{$\Delta\to\Delta^\wedge\to\Delta^\wedge\longrightarrow$}(\mathrm{Cat})\]
is a perfect test functor, where \hand{$\alpha:\Delta\to\Delta^\wedge$} is the canonical inclusion
It is tempting to surmise that this result is not special to i alone,
but that it holds for a large class, if not all, test functors from
\hand{$\Delta$} to asphericity or contractibility modelizers. Here Sd$^n$ denotes the
n.th iterate of the barycentric subdivision functor Sd (following
now the notation in Thomason's paper, which presumably is standard,
while I have been using ``Bar'' in part II of the notes). Presumably \inserted{functors} ana-%
loguous \deleted{functors} to Sd can be defined in any elementary modelizer A$^\wedge$,
as suggested by the natural constructions arising\unsure{\xout{,}} in connection with
the factorization property for a closed model structure on A$^\wedge$. Thus
possibly there is a general method in view for deducing perfect test
functors from ordinary one's. However, it definitely seems to me that
the natural place for these existence questions is in part V of the
notes, as they seem intimately related to a\hand{n} \deleted{closed} understanding of
the homotopy structures of \xout{categories} elementary modelizers, and more
\margin{$\uparrow$}specifically to the closed model structures to which such modelizers
\margin{$\lfloor$}give rise in various ways.

\hand{* * *}

\margin{$\downarrow$}Before proceeding, I would like to state still another after\xout{th}-%
thought to the reflections of section 89, about strictly generating
subcategories of a category M. I recall that a family of arrows in M
with same target X
\[u_i : X_i \longrightarrow X\]
is called \uline{strictly} \uline{epimorphic} , if for every object Y of M, the cor-%
responding map
\[(*)\qquad \mathrm{Hom}(X,Y)\longrightarrow \hand{$\prod_{i}$}\mathrm{Hom}(X_i,Y)\]
is injective (which is expressed by saying that the family $(u_i)$ is
\uline{epimorphic} ), and if \uline{moreover} the \xout{image of (*)} following obviously
necessary condition for an element $(f_i)$ of the product set of (*) to
be in the image of the map (*), is also sufficient:

(Comp) For any two indices i,j (possibly equal) and any commutat\unsure{i-}%
ve square
\diagram{\begin{tikzcd}[sep=small]
 & T \ar[dl, "\hand{$v_i$}"'] \ar[dr, "\hand{$v_j$}"] & \\
X_i \ar[dr, "\hand{$u_i$}"'] & & X_j \ar[dl, "\hand{$u_j$}"] \\
 & X &
\end{tikzcd}}
in M , the relation $f_iv_i = f_jv_j$ holds.

It is immediate that the condition for $(u_i)$ to be strictly epimorphic

\margin{123}
